\section{The fourth-order cancellation}
\label{subsec:interacting-fourth-jet}

Work in the finite ring
$\mathcal J=\C[t_c:c\in\mathcal C_*]/(t_c^5:c\in\mathcal C_*)$,
with coefficient norm
$\|P\|_r=\sum_{\alpha_c\le4}|P_\alpha|r^{|\alpha|}$.
This norm is submultiplicative, since discarding
terms with an exponent at least five can only
decrease the coefficient sum.
The jet map is a ring homomorphism; it
keeps degree at most four in each variable.
For $S=\{c\}$ set $k_S=\mathrm e^{w_c}-1$ in this
ring, and for $|S|\ge2$ set
$k_S=\mathrm e^{\mathcal G_S}-1$.
Every monomial of $k_S$ uses every variable
of $S$ and no other variable.

For a component $f_S$ bounded on a polydisc
of radius $R$, Cauchy's estimate at a smaller
radius $R'\le R/2$ gives
\[
 \|f_S\|_{R'}\le\sup|f_S|
             \left(\sum_{j=1}^4(R'/R)^j\right)^{|S|}
                         \le\sup|f_S|.
\]
Its monomials contain each variable of $S$
because it vanishes on every corresponding
coordinate hyperplane. Use this with
$R=t_x/2$, $R'=r$, and use the one-cell
estimate on radius $t_x$, to obtain
\begin{equation}\label{eq:interacting-error-jets}
 \|k_{\{c\}}\|_r\le\bar\kappa_1,\qquad
 \|k_S\|_r\le2\varpi_{\rm sep}^+
                       \le\bar\kappa_1^{N_*}\quad(|S|\ge2).
\end{equation}
We used $\|w_c\|_r\le3e_*^{1/4-9\varepsilon}\le3/8$
and $\|\mathrm e^P-1\|_r\le2\|P\|_r$ for $\|P\|_r\le1/2$.
The coefficients are continuous in $(\theta,A)$:
express the logarithm coefficients by Cauchy
integrals on a fixed small torus, and then
take the finite polynomial expressions in
those coefficients. They are real because
the logarithm is real at real marks.

In $\mathcal J$ the normalized generating
function has the exact representation
\begin{equation}\label{eq:interacting-jet-representation}
 \mathring\Phi=
   \prod_c\phi_0(t_c)
       \int\prod_{\varnothing\ne S\subset\mathcal C_*}
                         (1+k_S)\,v_{\rm sf}(\dd\theta\,\dd A)
       +L^{[4]}.
\end{equation}
To justify taking the jet inside the integral,
use joint continuity, holomorphy, and
$|\mathrm e^{\mathcal G}|\le2$ on a fixed
polydisc of radius at most $t_2$.
Cauchy's coefficient formulas and Fubini
then commute every required coefficient
with the finite measure integral.
Decompose $\mathcal G$ into its variable-set
components; its singleton components are
$\ell_0(t_c)+w_c(t_c)$.
Exponentiation, the jet homomorphism,
and $\mathrm e^{\ell_0}=N_0$ prove the representation.

\begin{theorem}\label{thm:interacting-reference-error}
Use the fixed construction, common trajectory and terminal inputs of
\cref{sec:rrm-terminal-inputs}. Suppose the factor and mass conditions
\eqref{eq:interacting-factor-conditions}--\eqref{eq:interacting-small-complex-mass},
the one-cell conditions \eqref{eq:interacting-one-cell-conditions},
the tree conditions \eqref{eq:interacting-tree-smallness}, and the jet
conditions \eqref{eq:interacting-jet-conditions} hold. Assume also
$r_*>31L+16$, $Mr_*/(8L)-1\ge6M+8$,
$v_*\ge8$, $p_0v_*^{p_0-1}\ge28$, and
$\lambda_L\le\bar\kappa_1^s$, where $0<s<1$.
For $m=2,4$,
\begin{equation}\label{eq:interacting-Theta-bound}
 \Theta_m^{(0)}\le
        m!r^{-m}\bigl(2\exp(1)\bar\kappa_1^s+\lambda_L\bigr)
                       \le8m!r^{-m}\bar\kappa_1^s.
\end{equation}
With
$\mathfrak e=s(1/4-9\varepsilon)-12\varepsilon>0$,
\begin{align}
 \Theta_4^{(0)}
    &\le3\cdot10^5 C_b^4e_*^{\mathfrak e}
           \le3\cdot10^5C_b^4h^{1/8}e^{\mathfrak e},\notag\\
 (1+|c_2|)\Theta_2^{(0)}
    &\le3\cdot10^4C_b^4e_*^{\mathfrak e}
           \le3\cdot10^4C_b^4h^{1/8}e^{\mathfrak e}.
                                    \label{eq:interacting-final-terminal-error}
\end{align}
The constants are independent of the ultraviolet cutoff and the
admitted placement. The terminal estimates used before harmonic
averaging hold uniformly in the real twist.
\end{theorem}

\begin{proof}
On the one-variable jet space define
$\ell_c(t_c^j)=m!\mathbf1_{\{j=m\}}-c_m\mathbf1_{\{j=0\}}$
and let $\ell=\bigotimes_c\ell_c$.
Then $(\Theta_m^{(0)})^{N_*}=|\ell(\mathring\Phi)|$.
The free jet satisfies $\ell_c(\phi_0)=0$.
From \eqref{eq:interacting-free-jet},
$|\ell(P)|\le(m!r^{-m})^{N_*}\|P\|_r$.
These statements follow by evaluating the
functionals on monomials with every exponent
in $\{0,m\}$.

Expand the product in
\eqref{eq:interacting-jet-representation} over
families $\mathcal Y$ of nonempty cell sets.
If $\mathcal Y$ fails to cover a cell $c$,
its error factors do not contain $t_c$;
the factor $\phi_0(t_c)$ is therefore
annihilated by $\ell_c$. Only covering
families contribute.
Assign cost one to each singleton and
cost $N_*$ to each larger set, and let
$a(\mathcal Y)$ be the sum of these costs.
A covering family has $a(\mathcal Y)\ge N_*$:
either it contains a larger set or it
contains every singleton.
By \eqref{eq:interacting-error-jets} its
coefficient norm is at most
$\bar\kappa_1^{a(\mathcal Y)}$.
The free factors have norm at most $2^{N_*}$.
Keep $\bar\kappa_1^{sN_*}$ from the cost
and sum the rest over all families. This
sum is
\[
 (1+\bar\kappa_1^{1-s})^{N_*}
  (1+\bar\kappa_1^{(1-s)N_*})^{2^{N_*}-1-N_*}
 \le\exp\left(N_*\bar\kappa_1^{1-s}
                  +2^{N_*}\bar\kappa_1^{(1-s)N_*}\right)
 \le\exp(1).
\]
The last step uses $N_*\bar\kappa_1^{1-s}\le1/2$.
The positive small-field measure has mass at
most one. Its total contribution is thus
bounded by
$(m!r^{-m})^{N_*}\exp(1)(2\bar\kappa_1^s)^{N_*}$.
The complement contributes at most
$(m!r^{-m}\lambda_L)^{N_*}$.
Taking the $N_*$th root proves
\eqref{eq:interacting-Theta-bound}.

Since $r^{-4}=256C_b^4e_*^{-12\varepsilon}$
and $\bar\kappa_1^s\le6e_*^{s(1/4-9\varepsilon)}$,
the $m=4$ bound has coefficient at most
$8\cdot24\cdot1536<3\cdot10^5$.
For $m=2$, $1+|c_2|\le(17/16)r^{-2}$,
so its coefficient is at most
$17\cdot1536<3\cdot10^4$.
Finally $0<\mathfrak e<1/4$ and
$e_*\le e\sqrt h$ imply
$e_*^{\mathfrak e}\le h^{1/8}e^{\mathfrak e}$.
The scale and probe bounds entering this
proof are uniform in the stated choices.
\end{proof}

